\section{Los datos como límite de la inteligencia: datos públicos, datos de dominio y datos generados por el producto}

Las dos secciones anteriores explican cómo se obtienen las capacidades por compresión; esta sección analiza qué determina el límite de esas capacidades. Wang Guan (王冠) reformula «tanta inteligencia como trabajo humano haya» como «tanta inteligencia como datos haya». Aquí los datos no son solo archivos o tablas: son la comprensión del problema, el know-how del negocio, el proceso de las tareas y la retroalimentación de los usuarios.

\begin{figure}[H]
\centering
\includegraphics[width=0.96\textwidth]{data-three-stages.png}
\caption{Las tres etapas de los datos: los datos públicos, los datos de dominio y los datos generados por el producto corresponden a distintos actores. Diagrama conceptual de elaboración propia, basado en la conversación de 00:37:11--01:05:36.}
\end{figure}

\begin{knowledgebox}{Lectura de la figura: cada etapa favorece a actores distintos}
La etapa de los datos públicos favorece a las empresas de modelos base, porque todas limpian y comprimen datos de internet; la etapa de los datos de dominio favorece a las grandes empresas tecnológicas y a las industrias tradicionales, porque poseen escenarios de negocio y datos privados; solo la etapa de los datos generados por el producto se parece más a una oportunidad para emprender con aplicaciones, porque un producto nuevo puede generar, al usarse, datos que antes no existían.
\end{knowledgebox}

\subsection{Una visión espacio-temporal de la inteligencia: datos, cómputo y algoritmos}

Wang Guan explica el límite de la inteligencia con una analogía en un espacio bidimensional. Los datos determinan la frontera del círculo, es decir, el alcance de capacidades al que el modelo puede llegar finalmente; el cómputo determina la velocidad con que se aproxima a esa frontera; los algoritmos determinan la trayectoria y la eficiencia de esa aproximación. Esta postura no niega la importancia de los algoritmos, sino que los sitúa dentro del límite que fijan los datos: sin los datos correspondientes, por potente que sea el algoritmo, difícilmente se obtienen capacidades estables de la nada.

\begin{importantbox}{Los datos determinan el límite de la inteligencia}
Si una tarea requiere know-how del sector, escenarios reales y juicio ante problemas abiertos, entonces «conseguir más datos» normalmente no es trabajo menor, sino que equivale a definir el problema mismo. Los gerentes de producto son importantes porque pueden convertir la comprensión del negocio en datos de los que se puede aprender.
\end{importantbox}

\subsection{Por qué las dos primeras etapas no son oportunidades típicas para emprender con aplicaciones}

Esta sección relaciona además las tres etapas de los datos con las oportunidades de emprendimiento. No todo tipo de datos está al alcance de una empresa emergente, ni toda etapa es adecuada para que entre una empresa de aplicaciones.

En la etapa de los datos públicos se compite en densidad de talento, cómputo, eficiencia de ingeniería y desgaste organizacional, por lo que resulta más adecuada para las empresas de modelos base. La etapa de los datos de dominio favorece a las grandes empresas tecnológicas y a las industrias con un alto grado de digitalización, porque ya cuentan con escenarios de negocio, canales de clientes y datos privados. La verdadera oportunidad para las empresas emergentes de aplicaciones está en la tercera etapa: los datos generados por el producto, es decir, los datos que surgen mientras el producto funciona, que antes no existían y que no pueden comprarse directamente fuera.

\begin{knowledgebox}{Ejemplos de datos generados por el producto}
\begin{tabularx}{\textwidth}{p{0.22\textwidth}p{0.34\textwidth}X}
\toprule
Tipo de datos & Cómo se generan & Por qué constituyen una barrera de entrada \\
\midrule
Trayectorias de intención del usuario & El usuario describe, corrige y elige continuamente dentro del producto & Reflejan necesidades y preferencias reales. \\
Muestras de generación fallida & El proceso de modificación cuando el contenido generado no satisface & Exponen las debilidades del modelo y los escenarios del producto. \\
Recipe & Métodos de producción y combinaciones de parámetros reutilizables & No son contenido aislado, sino capacidad de producción. \\
Retroalimentación del ciclo comercial & Si el contenido se reproduce, se paga, convierte o se reutiliza & Se vincula directamente con el valor del producto. \\
\bottomrule
\end{tabularx}
\end{knowledgebox}

\subsection{Resumen de la sección}

Las tres etapas de los datos explican por qué las empresas de aplicaciones no pueden limitarse a una capa delgada de UI. La verdadera oportunidad para las aplicaciones está en crear datos nuevos: cuanto más se usa el producto, más intenciones, procesos, fallas, recipes y retroalimentación comercial genera que no existen fuera de él.
